\documentclass[10pt, compress]{beamer}

\usetheme{glasgow}

\usepackage{booktabs}
\usepackage[scale=2]{ccicons}
\usepackage{minted}
\usepackage{bookmark}
\usepgfplotslibrary{dateplot}

\usemintedstyle{trac}

% ($ (A)!r!(B) $) the location of images to be used
\graphicspath{{src/}}

%% Customisation
% \newcommand{\V}[1]{\v} % vectors \v{c}
% \renewcommand{\v}[1]{\mathbf{#1}} % vectors
\newcommand{\ti}[1]{\tilde{#1}} % spectral representation
\newcommand{\tnsr}[1]{\underline{\underline{#1}}}

% Symbols
\renewcommand{\O}{\omega}  % omega
\newcommand{\E}{\varepsilon}  % epsilon
\renewcommand{\u}{\mu}  % mu
\newcommand{\p}{\rho}  % rho
\newcommand{\x}{\times}  % times
\renewcommand{\inf}{\infty}  % infinity
\newcommand{\infint}{\int\limits_{-\inf}^\inf} % integral by R
\newcommand{\e}{\mathrm{e}} % Straight-up exponential
\renewcommand{\j}{{j}\mkern1mu} % Straight-up exponential
\newcommand{\iu}{\mathrm{i}\mkern1mu}

\newcommand\ddfrac[2]{\frac{\displaystyle #1}{\displaystyle #2}}

\title{High Frequency Communication Systems}
\subtitle{Lecture 4}
\date{Spring 2021}
\author{Hasan T Abbas \& Qammer H Abbasi}
% \institute{}

\begin{document}

\maketitle

%%%%%%%%%%%%%%%%%%%%%%%%%%%%%%%%%%%%%%%%%%
%%%%%%%%%%%%%%%%%%%%%%%%%%%%%%%%%%%%%%%%%%
%%%%%%%%%%%%%%%%%%%%%%%%%%%%%%%%%%%%%%%%%%
\begin{frame}[fragile]
  \frametitle{Lecture Outline}
\begin{outline}[itemize]
  \1 Transmission Lines Theory \& Analysis
  \1 The Smith Chart 
  \1 Load Mismatching
  \1 The Magic of Quarter-wave Transformer
\end{outline}
\end{frame}
%%%%%%%%%%%%%%%%%%%%%%%%%%%%%%%%%%%%%%%%%%
%%%%%%%%%%%%%%%%%%%%%%%%%%%%%%%%%%%%%%%%%%
%%%%%%%%%%%%%%%%%%%%%%%%%%%%%%%%%%%%%%%%%%
\begin{frame}[fragile]
\frametitle{Transmission Line Theory}
\begin{outline}
  \1 A technique to bridge the gap between basic circuit analysis and electromagnetic fields theory
   \2 Lots of similaritites and analogies
    \1 Commonly used to design microwave devices and circuits
\end{outline}
\end{frame}

\begin{frame}
  \frametitle{TL \& EM Field Theory - Duality}
  \begin{columns}[T] % align columns
  \begin{column}{.3\textwidth}
\begin{outline}
  \1 In wave scattering problems, we place the co-ordinate system at the boundary
  \1 In transmission lines, we use the two different lines as the boundary
  \1 For open problems, the load is considered to be at infinity.
  \end{outline}
   \end{column}
 \begin{column}[T]{.7\textwidth}
    \begin{figure}
      \centering
          \includegraphics[width=.9\textwidth]{circuit.pdf}
      \caption{The analogy between wave problems and transmission lines.}
    \end{figure}
      \end{column}%
\end{columns}
  \end{frame}


\begin{frame}
  \frametitle{TL \& EM Field Theory Duality}
  \begin{table}
    \begin{tabular}{ll}
    \toprule
    TL Theory & EM Field Theory \\
    \midrule
    $V_{1} = V_{0} \e^{-\j \gamma_{1} z} \left(1+\Gamma_{L} \e^{\j 2 \gamma_{1} z}\right)$ & $E_{x,1}=E_{0} \e^{-\j k_{1} z} \left( 1 + \Gamma \e^{\j 2 k_{1} z} \right)$ \\
    $I_{1} = \frac{V_{0}}{Z_{0,1}} \e^{-\j \gamma_{1} z} \left(1+\Gamma_{L} \e^{\j 2 \gamma_{1} z}\right)$ & $H_{x,1} = \frac{E_{0}}{\eta_0} \e^{-\j k_{1} z} \left( 1 + \Gamma \e^{\j 2 k_{1} z} \right)$ \\
   $ \Gamma_{L} = \frac{Z_L - Z_{0,1}}{Z_L + Z_{0,1}} = \frac{Z_{0,2} - Z_{0,1}}{Z_{0,2} + Z_{0,1}}$ & $\Gamma = \frac{\eta_2 - \eta_1}{\eta_2 + \eta_1}$ \\
   $V_2 = T V_0 \e^{-\j \gamma_{2} z}$ & $E_{x,2} = T E_0 \e^{-\j k_{2} z}$ \\
   $I_2 = T \frac{V_0}{Z_{0,2}} \e^{-\j \gamma_{2} z}$ & $H_{x,2} = T \frac{E_0}{\eta_2} \e^{-\j k_{2} z}$ \\
   $T = \frac{2 Z_{0,2}}{{Z_{0,2}} + {Z_{0,1}}}$ & $T = \frac{2 \eta_2}{\eta_2 + \eta_1}$ \\
   \bottomrule
  \end{tabular}
\end{table}
\end{frame}


\begin{frame}
  \frametitle{TL Theory and Circuit Elements}
\begin{outline}
  \1 For a transmission line, we use a distributed circuit approach
  \2 Energy stored in the magnetic field $\rightarrow L$
  \2 Energy stored in the electric field $\rightarrow C$
  \2 Conductive losses $\rightarrow R$
  \2 Dielectric losses $\rightarrow G$
  \1 All the circuit elements are expressed per unit length
  \1 We can express the voltage and current at any given point $z$ and time $t$.
\end{outline}
\end{frame}

\begin{frame}
  \frametitle{The Transmission Line Equation}

\begin{figure}
  \centering
  \includegraphics[width=.85\textwidth]{tlines.pdf}
  \end{figure}
  \begin{figure}
    \centering
    \includegraphics[width=.85\textwidth]{circuit2.pdf}
      \end{figure}
\end{frame}

\begin{frame}
  \frametitle{Solution to the Transmission Line Equation - The Telegraphers Equation}
\begin{outline}
  \1 Solving the electric circuit leads to a second order differential equation
  \1 Analogous to the wave equation
  \2 Hence the analogies between voltage $V$ and electric field $\va{E}$
\end{outline}
\begin{align*}
  \left\{ \pdv[2]{z} - \left[ (R + \j \O L)( G + \j \O C)  \right] \right\} V (z,t) = 0
\end{align*}
The solution is of the type:
\begin{align*}
  V(z) &= V_{+} \e^{-\gamma z} + V_{-} \e^{+\gamma z}
\end{align*}
where the propagation constant $\gamma = \sqrt{(R + \j \O L)( G + \j \O C)}$. For a lossless case $R = G = 0$. 
\end{frame}

\begin{frame}
  \frametitle{Equivalent Circuits - Terminated Load}
\begin{outline}
  \1 
\end{outline}

  

\end{frame}

\end{document}
